\documentclass[11pt]{article}
\usepackage[utf8]{inputenc}
\usepackage[margin=1in]{geometry}
\usepackage{amsmath,amssymb}
\usepackage{booktabs}
\usepackage{graphicx}
\usepackage{xcolor}
\usepackage[colorlinks=true,linkcolor=blue,citecolor=blue,urlcolor=blue]{hyperref}
\usepackage{authblk}

% ============================================================================
% PAPER 2 (benchmark / systematic study) -- "Datasets & Benchmarks"-style
% All numbers are real measurements from the accompanying code release.
% Verify author/affiliation and bibliographic details before submission.
% ============================================================================

% >>> Target venue: NeurIPS 2026 Datasets & Benchmarks Track <<<
\title{\textbf{The Transform Zoo: A Systematic Benchmark of Orthogonal\\
Preconditioning for Vision--Language--Action Model Quantization}}
\author[1]{Justin Williams}
\author[1]{Kishor Datta Gupta}
\author[1]{Roy George}
\author[1]{Mrinmoy Sarkar}
\affil[1]{Clark Atlanta University \quad \texttt{justinwilliamstech693@gmail.com}}
\date{}

\begin{document}
\maketitle

\begin{abstract}
Orthogonal ``incoherence'' rotations are now standard for low-bit LLM quantization, but the
literature fixes on one choice (Hadamard, or a learned rotation) and the question
\emph{``which transform, and why?''} is largely unanswered---especially for closed-loop
Vision--Language--Action (VLA) policies. We benchmark \textbf{16 linear transforms}---fast
structured (Hadamard, DCT, DST, Hartley, Walsh, SRHT, butterfly, Haar), generic
(random-orthogonal, Givens, Householder, permutation), and data-dependent
(PCA, ZCA, polar)---as activation/weight preconditioners for 4-bit quantization of an
OpenVLA-OFT policy on LIBERO. We contribute: (i) a \textbf{cheap proxy} that ranks transforms
in seconds by their post-transform INT4 activation error, validated against full closed-loop
evaluation; (ii) a \textbf{taxonomy} organized by a single principle---\emph{quantization
wants energy spread uniformly across channels}---that explains every result, including the
striking negative one that PCA is actively harmful (it \emph{concentrates} energy, inflating
the outlier ratio to $1.7\times10^5$); and (iii) an \textbf{open library and leaderboard} so
practitioners can stop re-deriving these comparisons. Our headline practical result: any
global orthogonal mixer achieves $\sim\!90\%$ closed-loop success at W4A4, and \textbf{the DCT
is the deployment pick}---it ties Hadamard ($89.8\%$, $99.5\%$ retention) while applying the
true transform on \emph{every} layer (no power-of-two fallback) at $O(n\log n)$ cost.
\end{abstract}

\section{Introduction}
Incoherence processing---inserting an orthogonal $Q$ so that $Wx=(WQ^\top)(Qx)$ while the
rotated operands become near-Gaussian---underpins modern 4-bit LLM inference
\cite{quarot,spinquant,quip}. Yet practitioners face a confusing menu of transforms with
little guidance on which to use, why, or how to choose cheaply. For VLA policies
\cite{openvla,openvlaoft}, where the metric is \emph{closed-loop} task success and each
evaluation costs hours of simulation, the cost of trial-and-error is especially high.

This paper is a benchmark. We evaluate 16 transforms under one harness, propose and validate
a fast screening proxy, distill a predictive taxonomy, and release the code. Our aim is that
the community can read the taxonomy, use the library, and \emph{not} re-run the zoo.

\section{Transforms and Categories}
We cast every transform as an invertible preconditioner $M$ with
$y=Wx=(WM^{-1})(Mx)$, quantizing the conditioned weight $WM^{-1}$ and activation $Mx$. For
orthogonal $M$, $M^{-1}=M^\top$ and the error norm is basis-invariant; ZCA is the lone
non-orthogonal entry ($M=C^{-1/2}$). Table~\ref{tab:cat} groups the 16 transforms.

\begin{table}[h]\centering
\caption{The transform zoo, by data-dependence and structural cost.}
\label{tab:cat}
\begin{tabular}{lll}
\toprule
 & Fast $O(n\log n)$ structured & Dense $O(n^2)$ \\
\midrule
Data-free & Hadamard, DCT, DST, Hartley, Walsh, SRHT, butterfly, Haar & random-orth, Givens, Householder, permutation \\
Data-dependent & --- & PCA, ZCA, polar \\
\bottomrule
\end{tabular}
\end{table}

\section{A Cheap Proxy for Quantizability}
Closed-loop evaluation is slow; the mechanism is not. The failure mode of low-bit activation
quantization is outlier channels eating the per-tensor scale, so we score each transform by
the \emph{median INT4 per-tensor activation error} over backbone linears after applying $M$,
\begin{equation}
\mathrm{A4err}(x) \;=\; \frac{\mathbb{E}\,|\,\hat x - x\,|}{\mathbb{E}\,|x|},\qquad
\hat x = \mathrm{dequant}_{4}^{\text{per-tensor}}(Mx),
\end{equation}
computed offline from one calibration pass ($32$ frames). This ranks all 16 transforms in
seconds rather than the hours a rollout sweep would cost.

\section{Results}

\subsection{The ranking and the taxonomy}
Table~\ref{tab:rank} ranks all 16 transforms by the proxy. A single principle organizes the
outcome: \emph{quantization wants energy spread uniformly across channels}.

\begin{table}[h]\centering
\caption{Transform ranking by median post-transform INT4 activation error (lower is better;
identity $=96.3\%$). Outlier = median post-transform outlier ratio.}
\label{tab:rank}
\begin{tabular}{rlccl}
\toprule
\# & Transform & A4 err & Outlier $\rightarrow$ & Category \\
\midrule
1 & SRHT & 26.3\% & 1.74 & global mixer \\
2 & DST & 27.0\% & 1.83 & global mixer \\
3 & \textbf{DCT} & 27.3\% & 1.82 & global mixer (any-dim) \\
4 & Hadamard & 27.4\% & 1.80 & global mixer \\
5 & Walsh & 27.4\% & 1.80 & global mixer \\
6 & Hartley & 27.5\% & 1.86 & global mixer \\
7 & random-orthogonal & 28.2\% & 2.09 & global mixer \\
8 & polar & 70.7\% & 6.35 & partial \\
9 & butterfly & 71.0\% & 5.16 & partial \\
10 & PCA & 81.7\% & \textbf{171{,}211} & harmful (concentrates) \\
11 & Givens & 89.9\% & 12.17 & partial \\
12 & Haar & 90.9\% & 11.89 & localized \\
13 & identity & 96.3\% & 22.08 & control (no mixing) \\
14 & permutation & 96.3\% & 22.08 & control (reorders only) \\
15 & Householder & 96.4\% & 21.84 & too few reflections \\
16 & ZCA & 96.6\% & 8.52 & non-orthogonal \\
\bottomrule
\end{tabular}
\end{table}

\paragraph{Why each category lands where it does.}
\textbf{Global dense mixers} (ranks 1--7) blend every channel into every output, spreading
outliers $\Rightarrow$ work ($\sim\!27\%$). \textbf{Permutation} merely reorders ($=$ identity)
$\Rightarrow$ no help. \textbf{Haar} wavelets are \emph{localized} (each basis vector touches
few coordinates) $\Rightarrow$ cannot spread global outliers. \textbf{Householder ($k\!=\!24$)
/ Givens} are \emph{partial} rotations---a tunable ``amount-of-mixing'' knob.
\textbf{PCA} \emph{concentrates} energy into top components (its purpose for compression),
manufacturing a giant outlier---the exact opposite of what quantization needs.
\textbf{ZCA} (non-orthogonal whitening) amplifies the weight side and yields no net gain.

\subsection{Closed-loop validation of the proxy}
We then run closed-loop LIBERO on representatives of each tier (Table~\ref{tab:val}). The
proxy \textbf{correctly predicts ordering} (global mixers $\geq$ partial $>$ harmful), but is
\emph{pessimistic in magnitude}: it uses a per-tensor scale, whereas deployment uses
per-token quantization, which is more forgiving (PCA: predicted near-$0\%$, observed $50\%$).
We recommend the proxy for triage and closed-loop confirmation for finalists.

\begin{table}[h]\centering
\caption{Proxy vs.\ closed-loop W4A4 success. Screen = \textsc{spatial} ($30$ rollouts);
full = all suites ($400$ rollouts).}
\label{tab:val}
\begin{tabular}{lcccc}
\toprule
Transform & Proxy tier & Screen & Full eval & Predicted? \\
\midrule
random-orthogonal & global & 83.3\% & 90.5\% & \checkmark \\
DCT & global & 76.7\% & 89.8\% & \checkmark \\
SRHT & global & --- & 89.8\% & \checkmark \\
Hadamard & global & 80.0\% & 88.8\% & \checkmark \\
butterfly & partial & 73.3\% & --- & $\sim$ (ordering) \\
PCA & harmful & 50.0\% & --- & \checkmark (worst) \\
\bottomrule
\end{tabular}
\end{table}

\subsection{The deployment recommendation: DCT}
Among the tied global mixers, the DCT is the practical winner. Unlike Hadamard/SRHT---which
must fall back to a stored random matrix on the $32$ non-power-of-two
($11{,}008$-dim) \texttt{down\_proj} layers---the DCT applies the \emph{true} transform on all
$224$ backbone linears, is data-free, $O(n\log n)$, and full-eval-confirmed at $89.8\%$
($99.5\%$ retention). For any-dimension, training-free deployment, DCT dominates.

\section{The Library}
We release: (i) \texttt{rotations.py}, a registry of all 16 transforms with a uniform
\texttt{build\_transform(n, kind, ...)} API and exact-factorization unit tests; (ii) the
proxy driver \texttt{11\_transform\_compare.py}; and (iii) a persistent-model closed-loop
sweep harness with a leaderboard. New transforms are a few lines; new policies reuse the
harness.

\section{Limitations}
Simulated quantization (accuracy, not on-device latency); one VLA and one benchmark; the
proxy predicts ordering, not absolute magnitude (a per-token variant would tighten this).

\section{Conclusion}
Across 16 transforms, one principle predicts everything: spread energy, do not concentrate
it. Global orthogonal mixers tie at $\sim\!90\%$ W4A4 success; DCT is the deployment pick;
PCA is a cautionary tale; and a seconds-long proxy reproduces the ranking. We release the
benchmark so the community need not re-run it.

\begin{thebibliography}{99}
\small
\bibitem{openvla} A.~Kim et al. ``OpenVLA: An Open-Source Vision-Language-Action Model.'' 2024.
\bibitem{openvlaoft} M.~Kim et al. ``Fine-Tuning Vision-Language-Action Models (OpenVLA-OFT).'' 2025.
\bibitem{quarot} S.~Ashkboos et al. ``QuaRot: Outlier-Free 4-Bit Inference in Rotated LLMs.'' 2024.
\bibitem{spinquant} Z.~Liu et al. ``SpinQuant: LLM Quantization with Learned Rotations.'' 2024.
\bibitem{quip} J.~Chee et al. ``QuIP / QuIP\#: 2-Bit Quantization with Incoherence Processing.'' 2023--2024.
\end{thebibliography}

\end{document}
